\section{Conclusion}
\label{sec:conclusion}

Context-adaptive access control is useful precisely because it can lower its
guard, and that is also why it is dangerous. This paper argued that the
adaptation should be confined to a set that is already safe. We formalised a
constrained control problem over five latent threat states, four context
modalities and a $243$-element action lattice, with four invariants: an
unconditional actuation envelope (I1), an unconditional safety-traffic envelope
(I2), a belief-conditional restriction that can only narrow the set (I3), and
per-operation failure semantics that fail closed for confidentiality and
integrity operations and fail-safe-allow for certified safety operations (I4).
A forward Bayesian belief with cross-modal trust weights and a cross-node
dispersion signal chooses among the admissible actions; it has no authority over
the set itself.

The evaluation supports the separation. Across $800$ runs, every
invariant-enforcing policy records zero invariant violations per $1{,}000$
decisions, including in the \texttt{tee} scenario where the belief is confidently
wrong; the unconstrained risk-adaptive policy records $327.41$ and the static-low
policy $393.31$. The full controller reaches $0.120$ mean attack-scenario utility
against $0.084$ for static-high and $0.077$ for threshold CAAC, and denies
$0.029$ of safety-critical requests against $0.088$ for static-high. It is
cheaper than static-high in decision latency ($2.27$\,ms versus $6.00$\,ms) and
modelled energy ($0.288$\,mJ versus $0.728$\,mJ), and it reduces policy switching
from $15.3$ to $6.5$ per $1{,}000$ decisions relative to the same controller
without hysteresis.

The honest negative result is residual exposure. The invariant floor bounds the
mechanism of a granted action, not the credential of the requester: the full
controller grants $0.177$ of malicious requests against $0.127$ for static-high,
and the gap is concentrated in the out-of-model case where the node's trusted
execution environment is subverted ($0.473$ malicious grants, Brier $0.437$,
poisoning AUROC $0.273$). Within the stated threat model, the attested device
anchor and cross-node dispersion survive poisoning: the in-model downgrade yields
a Brier score of $0.003$ and a belief AUROC of $1.000$. Strengthening the
actuation floor to attestation plus deep inspection reduces malicious grants
without changing the zero-violation result, which is the right way to buy
exposure down once the estimator has been shown to be spoofable. The guarantee is
carried by the unconditional invariants; the belief improves utility and detection
when it can, and its failure does not become a safety failure.

Several directions remain. The most important is a hardware testbed that replaces
the analytical cost model with measured latency and energy on the three tiers. A
second is a stronger adversary that computes a policy against the defender's model
rather than following the scripted adaptive and resonant patterns; the invariant
argument would survive, but the utility and exposure numbers would be a better
bound. A third is out-of-band evidence---physical invariants, attestation rooted
outside the node TEE, or independent channels---that a full-modality spoof cannot
fabricate, which is the only way to close the \texttt{tee} gap rather than merely
bound its damage. Finally, the admissible-set constructor and the fail-safe
fallback are small enough to be verified with the timed-automata and
monitor-generation techniques that runtime assurance has developed, and doing so
would turn the structural argument of \cref{prop:structural} into a machine-checked
property of the implementation.

\section*{Data and code availability}

The simulator, the analysis scripts, the frozen $800$-run matrix and a
one-command reproduction procedure are available at
\url{https://github.com/DrMahmoudAljawarneh/failsafe-caci-iot}. The IoT-23 capture
used for the prior calibration is publicly available from the Stratosphere
Laboratory \cite{garcia2020iot}.
